% Part: model-theory
% Chapter: models-of-arithmetic
% Section: non-standard-models

\documentclass[../../../include/open-logic-section]{subfiles}

\begin{document}

% SEGMENT OLP-0194-S01

\olfileid{mod}{mar}{nst}
\section{நிலையற்ற மாதிரிகள்}

\begin{explain}
$\Lang{L_A}$ க்கான !!a{structure},~$\Struct{N}$ க்கு ஓரினமானது
எனில் அதை நிலையானது என்று அழைக்கிறோம். !!a{structure} ஒன்று
~$\Struct{N}$ க்கு ஓரினமல்ல எனில் அது நிலையற்றது எனப்படுகிறது.
\end{explain}

\begin{defn}
$\Lang{L_A}$ க்கான !!{structure}~$\Struct{M}$,~$\Struct{N}$ க்கு
ஓரினமல்ல எனில் \emph{நிலையற்றது}. ஏதேனும் $n \in
\Nat$ க்கு $\Value{\num{n}}{M}$ உடன் சமமான !!{element}s
$x \in \Domain{M}$, $\Struct{M}$ இன் \emph{நிலையான எண்கள்}
எனப்படுகின்றன; அவ்வாறு இல்லாதவை \emph{நிலையற்ற எண்கள்}.
\end{defn}

\begin{explain}
\olref[stm]{prop:standard-domain} படி,~$\Lang{L_A}$ க்கான எந்த
நிலையான !!{structure} உம் நிலையான !!{element}s ஐ மட்டுமே கொண்டுள்ளது.
எனவே, நிலையற்ற !!{structure} ஒன்று குறைந்தபட்சம் ஒரு நிலையற்ற
உறுப்பைக் கொண்டிருக்க வேண்டும். உண்மையில், நிலையற்ற !!{element}
ஒன்று இருப்பதே !!{structure} நிலையற்றது என்பதை உறுதிசெய்கிறது.
\end{explain}

\begin{prop}
$\Lang{L_A}$ க்கான !!a{structure}~$\Struct{M}$ ஒரு நிலையற்ற எண்ணைக்
கொண்டிருந்தால், $\Struct{M}$ நிலையற்றது.
\end{prop}

\begin{proof}
அவ்வாறில்லை என்று கொள்க; அதாவது $\Struct{M}$ நிலையானதாக இருந்தும்
நிலையற்ற எண்~$x$ ஒன்றைக் கொண்டுள்ளதாகக் கொள்க.
$g\colon \Nat \to \Domain{M}$ ஓர் ஓரினச்சார்பாகட்டும். $n$ மீது
தொகுத்தறிவதால் $g(\Value{\num{n}}{N}) =
\Value{\num{n}}{M}$ என்பதை எளிதாகக் காணலாம். வேறுவிதமாகச் சொன்னால்,
$g$ ஆனது~$\Struct{N}$ இன் நிலையான எண்களை~$\Struct{M}$ இன் நிலையான
எண்களுக்கு எடுத்துச்செல்கிறது. $\Struct{M}$ ஒரு நிலையற்ற எண்ணைக்
கொண்டிருந்தால், கருதுகோளுக்கு மாறாக $g$ !!{surjective} ஆக முடியாது.
\end{proof}

\begin{prob}
$\Th{Q}$ பின்வரும் அடிகோள்களைக் கொண்டிருப்பதை நினைவுகூர்க:
\begin{align*}
& \lforall[x][\lforall[y][(\eq[x'][y'] \lif \eq[x][y])]] \tag{$!Q_1$}\\
& \lforall[x][\eq/[\Obj 0][x']] \tag{$!Q_2$}\\
& \lforall[x][(\eq[x][\Obj 0] \lor \lexists[y][\eq[x][y']])] \tag{$!Q_3$}
\end{align*}
பின்வருமாறு அமையும் !!{structure}s~$\Struct{M_1}$, $\Struct{M_2}$,
$\Struct{M_3}$ ஆகியவற்றைத் தருக:
\begin{enumerate}
\item $\Sat{M_1}{!Q_1}$, $\Sat{M_1}{!Q_2}$, $\Sat/{M_1}{!Q_3}$;
\item $\Sat{M_2}{!Q_1}$, $\Sat/{M_2}{!Q_2}$, $\Sat{M_2}{!Q_3}$; மற்றும்
\item $\Sat/{M_3}{!Q_1}$, $\Sat{M_3}{!Q_2}$, $\Sat{M_3}{!Q_3}$;
\end{enumerate}
ஒவ்வொன்றுக்கும்~$\Assign{\Obj{0}}{M_i}$ மற்றும்
$\Assign{\prime}{M_i}$ ஐ மட்டும் குறிப்பிட வேண்டும் என்பது தெளிவு.
\end{prob}

\begin{explain}
$\Lang{L_A}$ க்கான நிலையற்ற !!{structure}s ஐக் குறிப்பிடுவது
எளிதானதே. எடுத்துக்காட்டாக, !!{domain}~$\Int$ உடைய கட்டமைப்பை எடுத்துக்
கொண்டு, தருக்கமல்லாக் குறியீடுகள் அனைத்தையும் வழக்கம்போல் பொருள்கொள்க.
எந்த~$n$ க்கும் எதிர்ம எண்கள் $\num{n}$ இன் மதிப்புகள் அல்ல என்பதால்,
இக்கட்டமைப்பு நிலையற்றது. உறுதியாக, இது~$\Struct{N}$ மெய்யாக்கும்
அதே வாக்கியங்களை மெய்யாக்கும் என்ற பொருளில் எண்கணிதத்தின்
\emph{மாதிரி} அல்ல. எடுத்துக்காட்டாக,
$\lforall[x][\eq/[x'][\Obj{0}]]$ பொய். எனினும், இறுதிக்குட்பட்ட
தன்மைத் தேற்றத்தைப் பயன்படுத்தி எண்கணிதத்தின் நிலையற்ற மாதிரிகள்
இருப்பதை எளிதாக நிரூபிக்கலாம்.
\end{explain}

\begin{prop}
$\Th{TA} = \Setabs{!A}{\Sat{N}{!A}}$ ஆனது~$\Struct{N}$ இன்
கோட்பாடாகட்டும். $\Th{TA}$ ஓர் !!a{enumerable} நிலையற்ற மாதிரியைக்
கொண்டுள்ளது.
\end{prop}

\begin{proof}
புதிய !!{constant}~$c$ ஒன்றால் $\Lang{L_A}$ ஐ விரிவாக்கி,
பின்வரும் !!{sentence}s இன் கணத்தைக் கருதுக:
\[
\Gamma = \Th{TA} \cup \{\eq/[c][\num{0}], \eq/[c][\num{1}],
\eq/[c][\num{2}], \dots\}
\]
$\Gamma$ இன் எந்த மாதிரி~$\Struct{M^c}$ உம் நிலையற்ற
!!a{element}~$x =
\Assign{c}{M}$ ஒன்றைக் கொண்டிருக்கும்; ஏனெனில் எல்லா
$n \in \Nat$ க்கும் $x \neq
\Value{\num{n}}{M}$. மேலும் $\Th{TA} \subseteq \Gamma$ என்பதால்
வெளிப்படையாக $\Sat{M^c}{\Th{TA}}$.~$c$ ஐ மறப்பதன் மூலம் மட்டும்
$\Struct{M^c}$ ஐ $\Lang{L_A}$ க்கான !!a{structure}~$\Struct{M}$
ஆக மாற்றினால், அதன் ஆட்களம் இன்னும் நிலையற்ற~$x$ ஐக் கொண்டிருக்கும்;
மேலும்~$\Sat{M}{\Th{TA}}$. $\Th{TA}$ இல் $c$ இடம்பெறாததால் பிந்தையது
உறுதிசெய்யப்படுகிறது. எனவே $\Gamma$ ஒரு மாதிரியைக் கொண்டுள்ளது என்று
காட்டினால் போதும்.

~$\Gamma$ ஒரு மாதிரியைக் கொண்டுள்ளது என்று காட்ட இறுதிக்குட்பட்ட
தன்மைத் தேற்றத்தைப் பயன்படுத்துகிறோம். ~$\Gamma$ இன் ஒவ்வொரு
இறுதிக்குட்பட்ட உட்கணமும் நிறைவுறத்தக்கது எனில்,~$\Gamma$ உம்
நிறைவுறத்தக்கது. ஏதேனும் இறுதிக்குட்பட்ட உட்கணம்
$\Gamma_0 \subseteq
\Gamma$ ஐக் கருதுக. $\Gamma_0$ ஆனது~$\Th{TA}$ இன் சில
!!{sentence}s ஐயும்~$\eq/[c][\num{n}]$ வடிவிலான சிலவற்றையும்
கொண்டுள்ளது; ஆனால் இறுதிக்குட்பட்ட எண்ணிக்கையில் மட்டுமே. இப்போது
$\eq/[c][\num{k}] \in \Gamma_0$ ஆகும் மிகப்பெரிய எண் $k$ என்று கொள்க.
பொருள்கோள்~$\Assign{c}{N_k} = k+1$ ஐச் சேர்க்கும் வகையில்
~$\Struct{N}$ ஐ விரிவாக்கி $\Struct{N_k}$ ஐ வரையறுக்க.
$\Sat{N_k}{\Gamma_0}$: $!A
\in \Th{TA}$ எனில், $c$ தவிர மற்றெல்லா வகையிலும்
~$\Struct{N_k}$ ஆனது~$\Struct{N}$ போலவே இருப்பதாலும்~$!A$ இல் $c$
இடம்பெறாததாலும், $\Sat{N_k}{!A}$. மேலும், $n \le k$ மற்றும்
$\Value{c}{N_k} = k+1$ என்பதால்,
$\Sat{N_k}{\eq/[c][\num{n}]}$. ஆகவே~$\Gamma$ இன் ஒவ்வொரு
இறுதிக்குட்பட்ட உட்கணமும் நிறைவுறத்தக்கது.
\end{proof}

\end{document}
